\section{Principal ideals in regular domains}

\begin{lemma}[Principal associated maximal ideals]\label{lem:regular-principal-associated-dvr}
\lean{EnvelopeProperty.isDiscreteValuationRing_of_regular_local_principal_associated}
\leanok
Let $(R,\mathfrak m)$ be a regular local domain. If $g\ne0$, $b\notin(g)$,
and $\mathfrak m b\subseteq(g)$, then $R$ is a discrete valuation ring.
\end{lemma}
\begin{proof}
Suppose $\dim R>1$. The finitely many primes minimal over $(g)$ have
height at most one. None contains $\mathfrak m$, and
$\mathfrak m\not\subseteq\mathfrak m^2$. Prime avoidance therefore gives
$x\in\mathfrak m\setminus\mathfrak m^2$ lying in none of these minimal
primes. The parameter quotient $R/(x)$ is regular local, hence a domain,
so $(x)$ is prime. If $g\in(x)$, the principal ideal height bound makes
$(x)$ minimal over $(g)$, contrary to the choice of $x$. Write
$xb=gc$. Primality of $(x)$ and $g\notin(x)$ give $c=xd$.
Cancellation of the nonzero $x$ gives $b=gd$, a contradiction. Thus the
embedding dimension of $R$ is at most one. A minimal generating family
of $\mathfrak m$ then has at most one member, so $\mathfrak m$ is
principal. The ring is not a field since $g\ne0$ and $b\notin(g)$.
A Noetherian local domain with nonzero principal maximal ideal is a DVR.
\end{proof}

\begin{lemma}[Associated primes of a principal quotient]\label{lem:regular-principal-associated-minimal}
\lean{EnvelopeProperty.isDiscreteValuationRing_localization_of_isAssociatedPrime_of_isRegularRing,EnvelopeProperty.mem_minimalPrimes_of_isAssociatedPrime_of_isRegularRing,EnvelopeProperty.associatedPrimes_eq_minimalPrimes_of_isRegularRing}
\leanok
\uses{lem:regular-principal-associated-dvr}
For a regular domain $E$ and $0\ne g\in E$, the associated primes of
$E/(g)$ are precisely the primes minimal over $(g)$, and all their
localizations are DVRs.
\end{lemma}
\begin{proof}
At an associated prime $P$, choose $b$ with $(g):b=P$. Its image in
$E_P$ is not in $(g)$: any contrary denominator clearing would give
$s\notin P$ with $sb\in(g)$. Every element of the maximal ideal of
$E_P$ times $b$ lies in $(g)$, so the preceding lemma makes $E_P$ a DVR.
For any prime $(g)\subseteq Q\subseteq P$, its extension to $E_P$ is a
nonzero prime, because localization of the domain is injective. It is
therefore the maximal ideal of this DVR. Contracting gives $Q=P$.
Thus each associated prime is minimal. Conversely, minimal primes over
the annihilator of a finite module over a Noetherian ring are associated;
the annihilator of $E/(g)$ is $(g)$.
\end{proof}

\begin{theorem}[DVR detection of principal ideals in regular domains]\label{thm:regular-principal-dvr-detection}
\lean{EnvelopeProperty.isDiscreteValuationRing_localization_of_mem_minimalPrimes_of_isRegularRing,EnvelopeProperty.mem_span_singleton_iff_forall_minimalPrimes_of_isRegularRing,EnvelopeProperty.mem_span_singleton_iff_forall_minimalPrimes_dvd_of_isRegularRing}
\leanok
\uses{lem:regular-principal-associated-minimal}
Let $E$ be a regular domain and $g\ne0$. For every prime $P$ minimal
over $(g)$, $E_P$ is a DVR. An element $f\in E$ belongs to $(g)$ if
and only if it belongs to $(g)E_P$ for every such $P$; equivalently,
for each such $P$ there is $s\notin P$ with $sf\in(g)$.
\end{theorem}
\begin{proof}
The DVR assertion follows from the associated-prime result. The forward
membership implication is immediate. If the class of $f$ in $E/(g)$
is nonzero, its annihilator is contained in an associated prime $P$.
This prime is minimal over $(g)$. An equality $sf\in(g)$ forces
$s$ into that annihilator and hence into $P$, contradicting $s\notin P$.
Finally, membership in the extended principal ideal is equivalent to
such a denominator-clearing equality by the definition of localization.
\end{proof}
